\section{Computational experiments}\label{sec:computational}

\subsection{Experimental design and evaluation protocol}\label{sec:compared}

The experiments compare how six algorithms solve sampled reduced-graph
problems and how the resulting firebreak placements perform on common
held-out propagation scenarios. We also specify a paired evaluation
for measuring transfer from reduced to full Reburn graphs.

The principal panel uses the 400-cell \texttt{new20x20} landscape and
its \texttt{shortest\_path} reduced propagation graphs. Its manifest
specifies training sizes $n\in\{100,200,400\}$, firebreak budget
fractions $\alpha\in\{0.01,0.02,0.03\}$, five independently generated
training samples per $(n,\alpha)$ configuration, and expected loss,
CVaR at confidence level $\beta=0.9$, and mean--CVaR with weight
$\lambda=0.5$. For each objective the same six configurations solve
each sampled problem, yielding $3\times3\times5\times3\times6=810$
recorded reduced-graph runs. The manifest specifies one thread, a
1,800\,s time limit, and requested relative MIP gap $10^{-3}$. The
preflight classifies 86.83\% of reduced scenarios as non-tree acyclic
DAGs, 12.41\% as rooted arborescences, and 0.76\% as empty or invalid.

\begin{center}
\begin{minipage}{0.96\linewidth}\centering\small
\captionof{table}{Six configurations matched within each sampled problem.
The three objectives use the corresponding exact labels in the manifest;
these are abbreviated consistently in the numerical tables.}
\label{tab:methods}
\begin{tabular}{p{0.37\linewidth}p{0.54\linewidth}}
\toprule
Displayed name & Configuration \\
\midrule
Extensive SAA & Extensive MIP with binary burn variables. \\
LP B\&B + Root & LP recourse cuts at integer candidates and LP-dual root cuts. \\
LP B\&B + Standard & Root cuts plus the initial standard LLBI. \\
LP B\&B + Coverage-exp & Root cuts and separated projected coverage inequalities. \\
LP B\&B + Path-exp & Root cuts and separated projected stored-path inequalities. \\
Combinatorial B\&B & Combinatorial candidate cuts with configured initial and fractional cuts. \\
\bottomrule
\end{tabular}
\end{minipage}
\end{center}

Within a block, methods share the objective, budget, training split and
fixed test split. Training IDs lie in 1--9000; all test splits contain
exactly IDs 9001--10000. The 810-run solver comparison is restricted
to reduced graphs. A separate Path-exp solution study uses the same
first five training splits per $(n,\alpha)$ and the fixed test set,
with five evaluated placements for each objective and design point.
Its solver certification is reported separately because the recorded
solution-study runs and the solver-comparison runs are distinct. The
reduced and Reburn \emph{training} splits differ and are not treated
as paired optimizations.

For any set $I$ of recorded runs, certification requires both the
\texttt{Optimal} status and reported final relative MIP gap
$g_i\leq 10^{-3}$. We report $100N_{\rm opt}/|I|$ (\%), mean and maximum final
relative MIP gaps in percent, and mean observed solver time
$|I|^{-1}\sum_{i\in I}t_i$. Every statistic includes unresolved runs;
positive gaps within the tolerance and recorded time overruns are
retained. Mean observed runtime is computation consumed under the
stopping rule, not an estimate of unrestricted time to certification:
unsuccessful runs are right censored at the stopping horizon. Our
cumulative certification profiles count certifications by observed
solver time over a common denominator of 45 per objective and method.
They are not Dolan--Mor\'e profiles, which would require a reliable
uncensored paired time-to-solution for every compared case.
